\section{Synthesizing Rewrites with Dataflow Guidance}
\label{slab:sec:synthesis}
\label{slab:sec:alg}

The abstraction of \cref{slab:sec:dataflow} becomes an optimizer only
inside a search loop: something must generate candidates, consult
their dataflows, discard the hopeless, verify the promising, and keep
the fastest. This section describes that machinery. Its core, the
scoring function of \cref{slab:sec:scorefunction} and the stratified
search of \cref{slab:sec:prioritizingsearch}, is the contribution of
Rui Dong; both are summarized here to keep the chapter self-contained,
and the reader is referred to the publication~\cite{dong2023slabcity}
for their complete development. The equivalence-checking and ranking
components (\cref{slab:sec:checkingequiv,slab:sec:perfranking}) were
developed collaboratively.

\subsection{The Top-Level Loop}
\label{slab:sec:toplevelalg}

Algorithm~\ref{slab:alg:toplevelalg} shows \toolname end to end. Given
$\inputquery$ and $\integrityconstraint$, it instantiates the CEGIS
paradigm~\cite{solar2008program} introduced in Chapter~2 for query
rewriting: an inductive synthesizer proposes candidates consistent
with the counterexamples accumulated so far, and the equivalence
checker either admits or rejects each candidate. New counterexamples
are produced by the checker's first two layers
(\cref{slab:sec:checkingequiv}): when the tester or the bounded
verifier refutes a candidate, the distinguishing database it found,
paired with the output $\inputquery$ produces on it, joins the
accumulated set (line~\ref{slab:alg:toplevelalg:addcounterexample}).
Checking a candidate against that set
(line~\ref{slab:alg:toplevelalg:satisfycounterexample}) is just query
execution, far cheaper than invoking a verifier
(line~\ref{slab:alg:toplevelalg:checkequiv}). This is the economy
CEGIS buys: most wrong candidates die on replayed counterexamples, and
the expensive checker is consulted only for candidates that survive
everything already known.

\begin{figure}[!t]
\vspace{-5pt}
\footnotesize
\centering

\begin{algorithm}[H]
\begin{algorithmic}[1]

\myprocedure{Optimize}{$\inputquery, \integrityconstraint$}

\Statex\Input{An input query $\inputquery$ and an integrity constraint $\integrityconstraint$.}


\Statex\Output{An equivalent and faster output query $\query$.}

\vspace{2pt}

\State $\counterexamples \assign \emptyset$; \  $\queryset \assign \emptyset$; 
\label{slab:alg:toplevelalg:init}



\While{not timed out}
\label{slab:alg:toplevelalg:whilebegin} 



\ForAll{$\query \in \textsc{LazyEnumerate}(\inputquery)$} \label{slab:alg:toplevelalg:forloop}
\Comment{{{\scriptsize \Circled{1}}} Query synthesizer.}


\State \textbf{if} $\query$ doesn't pass  $\counterexamples$ \textbf{then} \textbf{continue}; \label{slab:alg:toplevelalg:satisfycounterexample}
\Comment{Check against counterexs.}


\State $\counterexample \assign \textsc{CheckEquivalence}(\query, \inputquery, \integrityconstraint)$;\label{slab:alg:toplevelalg:checkequiv}
\Comment{{{\scriptsize \Circled{2}}} Equivalence checker.}

\vspace{1pt}

\If{$\counterexample = \emph{null}$} 
$\queryset.\emph{add}(\query)$
\label{slab:alg:toplevelalg:nocounterexample}
\Comment{$\emph{null}$ means $\query$ is equivalent to $\inputquery$.}

\Else \ $\counterexamples.\emph{add}( \counterexample )$;
\Comment{Otherwise, $\counterexample$ is a new counterexample.}
\label{slab:alg:toplevelalg:addcounterexample}

\EndIf

\EndFor


\EndWhile
\label{slab:alg:toplevelalg:whileend}

\vspace{-1pt}
\State $\query \assign \textsc{RankPerformance}(\queryset, \inputquery)$;
\label{slab:alg:toplevelalg:checkfaster}
\Comment{{{\scriptsize \Circled{3}}} Performance Ranker.}

\State \Return $\query$;
\label{slab:alg:toplevelalg:returnquery}

\caption{{Top-level algorithm of \toolname.}}
\label{slab:alg:toplevelalg}
\end{algorithmic}
\end{algorithm}
\vspace{-10pt}
\end{figure}

The synthesizer does not enumerate blindly. Candidates are generated
lazily (line~\ref{slab:alg:toplevelalg:forloop}) in non-ascending
order of their dataflow-based score, defined next, so that the queries
most likely to be equivalent to and faster than $\inputquery$ surface
first. Candidates the checker admits accumulate in a result set
(line~\ref{slab:alg:toplevelalg:nocounterexample}); when the time
budget (five seconds by default) expires, the performance ranker
(\cref{slab:sec:perfranking}) orders the result set
(line~\ref{slab:alg:toplevelalg:checkfaster}) and the best query is
returned, carrying the verification label
(\cref{slab:sec:checkingequiv}) it earned.

\subsection{Scoring Candidates by Dataflows}
\label{slab:sec:scorefunction}

The score realizes, as a single number, the one-sided signal of
\cref{slab:sec:dataflow:insight}. For a program $\program$ arising in
a context $\query_{\emph{ctx}}$, extraction
(\cref{slab:sec:dataflow:extraction}) yields $\program$'s dataflow
multiset $\dataflowset$, and the input query's multiset
$\dataflowset_{\emph{in}}$ is extracted once, up front, closed under
alternative dataflows. The score of $\program$ is the negated count of
$\program$'s dataflows that are foreign to the input:
\[
\scorefunction(\program \mid \query_{\emph{ctx}})
= - \, \big| \, \dataflowset \setminus \dataflowset_{\emph{in}} \, \big|.
\]
Zero is therefore the best possible score, attained exactly when every
computation $\program$ performs already occurs in $\inputquery$; each
foreign dataflow costs one point. Sharing is rewarded, foreignness is
penalized, and using less than everything $\inputquery$ computes costs
nothing, which is the asymmetry that lets optimized queries shed the
input's redundant work without being punished for it.

Because the score counts dataflows and dataflows ignore syntax, the
score inherits the semantic character of the abstraction: a large,
structurally ambitious candidate whose computations all come from
$\inputquery$ outscores a tiny candidate that introduces one foreign
computation. Ties between equal-scoring candidates are broken toward
syntactically smaller queries, so size ordering survives as a tiebreak
rather than as the organizing principle.

Finally, the compositionality property of
\cref{slab:sec:dataflow:extraction} transfers to the score directly:
since a component's dataflow multiset is contained in the whole
program's, a component's score is always at least the whole program's
score. Composing programs can only preserve or worsen the score, never
improve it. This monotonicity is what the search exploits.

\subsection{Stratified Search}
\label{slab:sec:prioritizingsearch}

Monotonicity turns the score into a pruning device: if every building
block available at some score level is known, then no composition of
them can score higher, so the space can be explored in strata. The
search algorithm enumerates queries layer by layer, first exhausting
all candidates with score $0$ without ever constructing anything
lower-scoring, then all candidates with score $-1$, and so on. Within
a stratum, construction is bottom-up and follows a dynamic-programming
design: to generate the queries of score $\score$, it suffices to
combine queries, conditions, and target lists whose scores are at
least $\score$, and the sets of those are exactly what the previous
strata computed. A set of syntax-directed inference rules governs
which compositions are attempted, mirroring the query grammar of
\cref{slab:sec:problem}.

One subtlety concerns termination within a stratum. If dataflows were
collected into sets, a generator could manufacture unboundedly many
new queries sharing one dataflow set, and a stratum would never close.
Representing dataflows as multisets removes the loophole: repeating a
computation grows the multiset, changes the difference against
$\dataflowset_{\emph{in}}$, and lowers the score, so each stratum is
finite.

This summary deliberately omits the search algorithm's full pseudocode
and its complete rule set, which are Rui Dong's contribution and are
presented in detail in the publication~\cite{dong2023slabcity}.

\subsection{Checking Query Equivalence}
\label{slab:sec:checkingequiv}

Every admitted candidate must be defended, and Chapter~2 laid out the
spectrum of available defenses: testing, bounded verification, full
verification, ordered by cost and by strength of guarantee. \slabcity
uses all three, cheapest first, and this layering is what makes the
checker affordable inside a search loop that proposes candidates by
the hundreds.

\head{Syntax-guided query testing}
The first layer is a tester built specifically for this system.
Off-the-shelf equivalence testing for a language as expressive as ours
barely exists, and randomly generated databases are almost useless
here: they overwhelmingly produce empty or degenerate outputs on which
inequivalent queries agree. Our tester instead mines the queries
themselves for hints about which databases can tell them apart.
\begin{itemize}[leftmargin=*]
\item
\emph{Satisfying filter and join conditions.}
It reads filter conditions (such as those in \sqlwherecolor) from
$\inputquery$ and populates test databases with rows satisfying them,
and likewise collects \sqljoincolor conditions, so that outputs and
intermediate join results are non-empty; non-empty outputs are what
disprove equivalence.
\item
\emph{Duplicating non-key values.}
It looks for operations whose misplacement is a common source of wrong
candidates, \sqldistinctcolor above all, and generates duplicated
values in the affected columns: if a candidate places \sqldistinctcolor
on a column $C$, a test database with duplicates in $C$ is exactly
what makes a misplaced \sqldistinctcolor change the output.
\item
\emph{Duplicating grouping columns.}
It compares grouping structure: when a candidate's \sqlgroupbycolor
columns differ from the input's, it duplicates values in exactly those
columns, the situation under which divergent groupings disagree.
\end{itemize}
Each family of hints is encoded as a logical formula $\varphi_1$,
conjoined with a formula $\varphi_2$ encoding the integrity constraint
$\integrityconstraint$, and handed to a constraint solver (such as
Microsoft Z3); each satisfying assignment of
$\varphi_1 \land \varphi_2$ is a test database. When the conjunction
is unsatisfiable, no database can realize the hints under the
constraints, and the candidate proceeds to verification.

\head{Bounded verification}
The second layer is bounded verification, reusing the bounded
verifiers of \cosette~\cite{chu2017cosette,wang2018speeding}. Its
exhaustiveness within the bounded space cuts both ways: a refutation
comes with a distinguishing database, which joins the counterexample
set exactly as the tester's do, while a positive verdict guarantees
that $\inputquery$ and $\query$ agree on every database with at most
$k$ rows per table ($k=2$ by default, customizable, with cell values
unbounded). The verifier may also reach no conclusion within its
budget, in which case the candidate retains only the tester's implicit
endorsement.

\head{Full verification}
The final layer is full verification via \cosette and
\spes~\cite{chu2017cosette,zhou2022spes}, which proves equivalence
over all possible inputs but succeeds mostly on simpler query shapes
(such as select-project-join) and costs the most; \slabcity invokes it
parsimoniously, only for candidates nothing else could refute, and its
refutations, unlike the earlier layers', are verdicts without
witnesses. Each surviving candidate thus carries one of three labels,
in decreasing strength: fully verified, bounded-verified, or tested;
\cref{slab:sec:eval:discussion} discusses what the labels mean for
deployment.

\subsection{Performance Ranking}
\label{slab:sec:perfranking}

Equivalence is only half of \cref{slab:def:problem}; the other half is
speed. The performance ranker orders the admitted candidates by
estimated cost from \texttt{EXPLAIN} and returns the winner. Cost
estimates are a proxy, and a deliberately cheap one; where the
overhead is justified, \slabcity can instead execute candidates on a
sample of the data, or on the full data for queries that will be rerun
indefinitely, such as reporting dashboards. How performance is
predicted is orthogonal to how candidates are found, and
\cref{slab:sec:ablation} quantifies how little the choice costs:
ranking by \texttt{EXPLAIN} comes close to ranking by true latency.
